
\documentclass[12pt]{article}
\usepackage[spanish]{babel}
\usepackage[utf8]{inputenc}
\usepackage[T1]{fontenc}
\usepackage{geometry}
\usepackage{graphicx}
\usepackage{booktabs}
\usepackage{hyperref}
\usepackage{longtable}
\geometry{margin=2.5cm}

\title{Informe Ejecutivo\\
\large Sistema de Realidad Aumentada e Inteligencia Artificial para la Gestión de Plagas y Enfermedades en el Cultivo de Arroz en el Meta, Colombia}
\author{Julio Vargas-Riaño}
\date{Febrero 2026}

\begin{document}

\maketitle
\tableofcontents
\newpage

\section{Resumen Ejecutivo}
Este informe presenta un análisis integral y comparativo (benchmarking) de modelos de Inteligencia Artificial (IA) y Visión por Computadora orientados a la detección automática de plagas y enfermedades en el cultivo de arroz, con énfasis en la viabilidad técnica para dispositivos móviles de gama de entrada y gafas de Realidad Aumentada (RA), bajo condiciones de operación \textbf{offline} en zonas rurales del departamento del Meta, Colombia.

\section{Objetivos}

\subsection{Objetivo General}
Desarrollar un sistema inteligente basado en IA y RA para la detección temprana, clasificación y gestión de plagas y enfermedades del arroz.

\subsection{Objetivos Específicos}
\begin{itemize}
\item Evaluar modelos de visión por computadora bajo restricciones de hardware.
\item Integrar modelos ligeros con sistemas RA.
\item Permitir diagnóstico en tiempo real sin conectividad.
\item Reducir el uso innecesario de pesticidas.
\end{itemize}

\section{Metodología}
Se realizó una revisión sistemática de literatura (2016–2025), análisis técnico de arquitecturas CNN y YOLO, evaluación experimental y análisis de viabilidad en hardware móvil.

\section{Benchmarking de Modelos}

\begin{longtable}{lcccc}
\toprule
Modelo & Arquitectura & Precisión & Parámetros & Offline \\
\midrule
YOLOv4-tiny & YOLO Tiny & 97.36\% mAP & 6.0M & Sí \\
YOLOv5s + Ghost & YOLO + Ghost & 96.1\% mAP & 3.9M & Sí \\
LiSA-MobileNetV2 & MobileNetV2 + Attn & 95.7\% Acc & 2.0M & Sí \\
GBiDC-PEST & YOLO + GhostNet & 80.1\% mAP & 1.2M & Sí \\
EfficientNet-B0 & EfficientNet & 96\% Acc & 5.3M & Sí \\
Simple CNN & CNN ligera & 93.3\% Acc & 0.8M & Sí \\
\bottomrule
\end{longtable}

\section{Análisis Técnico}

\subsection{Ventajas}
\begin{itemize}
\item Alta eficiencia computacional.
\item Operación offline.
\item Detección multiclase.
\item Integración con realidad aumentada.
\end{itemize}

\subsection{Limitaciones}
\begin{itemize}
\item Compromiso precisión-velocidad.
\item Sensibilidad a iluminación variable.
\item Escalabilidad multiclase limitada.
\end{itemize}

\section{Propuestas Técnicas}

\subsection{YOLOv4-tiny}
Modelo principal de detección por su latencia ultra baja y alta precisión.

\subsection{LiSA-MobileNetV2}
Clasificador optimizado para diagnóstico preciso con bajo consumo de memoria.

\subsection{CNN Ligera}
Modelo de respaldo para hardware extremadamente limitado.

\section{Arquitectura del Sistema}
Captura de imagen → Preprocesamiento → Detección → Clasificación → Visualización RA → Recomendaciones.

\section{Impacto Esperado}
\begin{itemize}
\item Reducción del 30–45\% en pesticidas.
\item Incremento del 10–25\% en rendimiento del cultivo.
\item Mejora ambiental y económica.
\end{itemize}

\section{Recomendaciones}
\begin{itemize}
\item Cuantización INT8.
\item Limitar a 3–6 enfermedades críticas.
\item Integrar interfaz de voz offline.
\item Capacitación rural.
\end{itemize}

\section{Fuentes}
\begin{enumerate}
\item Deng et al. (2023). Agronomy.
\item Rahman et al. (2018). Biosystems Engineering.
\item Khanal et al. (2024). arXiv.
\item Su et al. (2024). Journal of Integrative Agriculture.
\item Jain et al. (2022). Electronics.
\item Nayak et al. (2023). Smart Agricultural Technology.
\item Hasan et al. (2023). Agriculture.
\item Pan et al. (2024). Frontiers in Plant Science.
\end{enumerate}

\section{Conclusión}
La integración de IA ligera y RA representa una solución viable, escalable y sostenible para la agricultura de precisión en Colombia.

\end{document}
